% ============================================================================
% Fichier  : partie3/P03_C11_audit_controle_assurance.tex
% Titre    : Audit, Contrôle et Assurance de la Gouvernance des SI
% Auteur   : MINKA MI NGUIDJOÏ Thierry Emmanuel
% Version  : 1.0 -- octobre 2026
% ============================================================================
% Sources du dossier Normes : IIA, Cadre de référence international des
% pratiques professionnelles (édition 2017) et Code de déontologie ; INTOSAI
% GUID 5100 (2019) ; IDI, Manuel de mise en œuvre des ISSAI d'audit de
% conformité (2021). Les Normes mondiales de l'IIA (2024) sont présentées
% d'après leur structure publique. Laboratoire associé : GCI-LAB-05 (test des
% contrôles généraux informatiques à l'ANPP).
% ============================================================================

\chapter{Audit, Contrôle et Assurance de la Gouvernance des SI}
\label{chap:audit}
\chaptermark{Audit, contrôle et assurance}

\epigraph{%
    « Une réponse orale n'est pas une preuve. »%
}{Règle des auditeurs}

\niveauChapitre{Intermédiaire}{%
    Maîtrise des Chapitres~\ref{chap:cobit-capacite} (évaluation de
    capacité) et~\ref{chap:droit} (obligations légales).%
}

\begin{boxobjectifs}
\begin{itemize}
    \item Appliquer le modèle des trois lignes.
    \item Situer les Normes mondiales d'audit interne de l'IIA (2024) et le
          cadre de 2017.
    \item Mener un audit des SI selon l'INTOSAI GUID~5100 dans le secteur
          public.
    \item Tester des contrôles généraux informatiques (accès, changements,
          exploitation).
    \item Rédiger un constat d'audit qui sépare constatation, critère,
          cause, conséquence et recommandation.
    \item Identifier les exigences déontologiques de l'auditeur et du DSI.
\end{itemize}
\end{boxobjectifs}

\bigskip

% ============================================================================
\section*{Question directrice : qui vérifie que tout cela fonctionne~?}
% ============================================================================

Les chapitres précédents ont construit un système de gouvernance, fixé des
capacités cibles et dressé la liste des obligations légales. Il reste une
question~: \emph{qui vérifie, et avec quelle indépendance, que ce système
fonctionne réellement~?} À Yaoundé, la Chambre des comptes vient d'annoncer
à l'ANPP (opération MEFOU) un audit de ses systèmes d'information. À
Douala, Mme~EKAMBI, directrice de l'audit interne de la BANQUE SANAGA, doit
expliquer au comité d'audit pourquoi sa recommandation d'octobre~2025 est
restée sans suite.

Ce chapitre présente l'organisation de l'assurance (le modèle des trois
lignes, section~\ref{sec:c11-trois-lignes}), les normes de l'audit interne
(section~\ref{sec:c11-iia}) et de l'audit public
(section~\ref{sec:c11-intosai}), la pratique du test des contrôles
généraux informatiques (section~\ref{sec:c11-itgc}), la rédaction des
constats (section~\ref{sec:c11-constat}) et la déontologie
(section~\ref{sec:c11-deontologie}). Le laboratoire GCI-LAB-05 fait tester des
contrôles à partir des pièces de l'ANPP.

\separateur

% ============================================================================
\section{Le Modèle des Trois Lignes}
\label{sec:c11-trois-lignes}
% ============================================================================

\subsection{Qui fait quoi dans la maîtrise des risques}

L'Institut des auditeurs internes (\sigle{IIA}) a publié en 2020 le
\termecle{modèle des trois lignes}\index{modèle des trois lignes}
(\textit{Three Lines Model}), qui a succédé au modèle des « trois lignes de
défense ». Il répartit les rôles entre l'organe de gouvernance, le
management et l'audit interne~\cite{iia2020troislignes}.

\begin{tableaucours}{L{2.8cm}L{4.2cm}L{5.2cm}}{%
    \tableauHeader{Acteur} &
    \tableauHeader{Rôle} &
    \tableauHeader{À la BANQUE SANAGA}
}
Organe de gouvernance & Rend compte aux parties prenantes, fixe la
direction, supervise & Conseil d'administration, comité d'audit, comité
stratégique des SI \\
Première ligne & Le management opérationnel, qui porte et maîtrise les
risques dans ses activités & DSI, exploitation, chefs de projet, métiers \\
Deuxième ligne & Les fonctions d'expertise, de soutien et de surveillance
des risques & Gestion des risques, conformité, RSSI, responsable GRC \\
Troisième ligne & L'audit interne, qui fournit une assurance indépendante
et objective & Direction de l'audit interne (Mme~EKAMBI) \\
Prestataires d'assurance externes & Assurance complémentaire,
indépendante de l'organisation & Commissaire aux comptes, COBAC, audit de
sécurité prévu par la loi n°~2010/012 \\
\end{tableaucours}
\vspace{-0.8em}
\begin{center}{\small\itshape Tableau~11.1. Le modèle des trois lignes
appliqué à la BANQUE SANAGA}\end{center}

\subsection{Trois erreurs fréquentes}

\begin{enumerate}
    \item \textbf{Confier la deuxième ligne à la première.} Un RSSI
          rattaché au DSI, et noté par lui, ne peut pas surveiller librement
          les risques que le DSI prend. Le modèle ne prescrit pas un
          rattachement, mais il exige que la deuxième ligne puisse remonter
          ses alertes sans filtre.
    \item \textbf{Faire faire le travail de la première ligne par la
          troisième.} Un audit interne qui rédige les procédures qu'il
          auditera ensuite perd son indépendance.
    \item \textbf{Croire que l'audit est responsable du risque.} La
          troisième ligne donne une assurance~; elle ne maîtrise pas le
          risque. La recommandation non suivie de Mme~EKAMBI n'est pas une
          faute de l'audit~: c'est une défaillance du suivi par la première
          ligne et par le comité d'audit (MEA02).
\end{enumerate}

% ============================================================================
\section{Les Normes de l'Audit Interne}
\label{sec:c11-iia}
% ============================================================================

\subsection{Du cadre de 2017 aux Normes mondiales de 2024}

L'IIA définit l'audit interne comme une activité indépendante et objective
qui donne à une organisation une assurance sur le degré de maîtrise de ses
opérations, lui apporte des conseils pour les améliorer et contribue à
créer de la valeur, en évaluant de manière systématique ses processus de
gouvernance, de management des risques et de contrôle~\cite{iia2017cripp}.

Jusqu'en 2024, la profession se référait au \emph{Cadre de référence
international des pratiques professionnelles} (\sigle{CRIPP}), dans son
édition de 2017, disponible en français dans le dossier \texttt{Normes}.
Il comprenait notamment la définition de l'audit interne, le Code de
déontologie et des normes réparties entre normes de qualification (série
1000) et normes de fonctionnement (série 2000).

En janvier~2024, l'IIA a publié les \textbf{Normes mondiales d'audit
interne} (\textit{Global Internal Audit Standards}), applicables depuis
janvier~2025~\cite{iia2024gias}. Elles remplacent le CRIPP et
s'organisent en cinq domaines~:

\begin{tableaucours}{cL{4.0cm}L{8.0cm}}{%
    \tableauHeader{Dom.} &
    \tableauHeader{Intitulé} &
    \tableauHeader{Contenu}
}
I & Raison d'être de l'audit interne & Ce que l'audit interne apporte à
l'organisation \\
II & Éthique et professionnalisme & Intégrité, objectivité, compétence,
conscience professionnelle, confidentialité \\
III & Gouvernance de la fonction d'audit interne & Relation avec le
conseil, mandat, indépendance, supervision \\
IV & Management de la fonction d'audit interne & Stratégie, ressources,
méthodologie, qualité \\
V & Réalisation des missions & Planification, conduite, communication des
résultats, suivi \\
\end{tableaucours}
\vspace{-0.8em}
\begin{center}{\small\itshape Tableau~11.2. Les cinq domaines des Normes
mondiales d'audit interne (2024, notre traduction)}\end{center}

Le domaine~III est le plus directement lié à la gouvernance~: il place
expressément sous la responsabilité du conseil la supervision de la
fonction d'audit interne, son mandat et les moyens de son indépendance.

\begin{boxalert}
\textbf{Violation typique.} La charte d'audit interne de la BANQUE SANAGA
date de 2015 et renvoie aux « normes de l'IIA » sans préciser l'édition.
Le plan d'audit est approuvé par le directeur général, et non par le
comité d'audit.

\smallskip
\textbf{Règle~:} la charte d'audit doit se référer aux Normes mondiales de
2024, et le plan d'audit comme les moyens de la fonction doivent être
approuvés par l'organe de gouvernance, non par la direction auditée.
\end{boxalert}

% ============================================================================
\section{L'Audit Public des SI : l'INTOSAI GUID 5100}
\label{sec:c11-intosai}
% ============================================================================

\subsection{Les institutions supérieures de contrôle}

Dans le secteur public, l'audit externe relève des institutions
supérieures de contrôle des finances publiques (\sigle{ISC}), qui
appliquent les normes de l'INTOSAI (les \sigle{ISSAI}). Au Cameroun, cette
fonction est assurée par la Chambre des comptes de la Cour suprême. Le
GUID~5100 de l'INTOSAI (2019) fournit le cadre général de la vérification
des systèmes d'information~\cite{intosai2019guid5100}.

\begin{boxdef}
\termecle{Vérification des systèmes d'information}\index{vérification des
systèmes d'information} (\textit{information systems audit})~: selon le
GUID~5100, examen des contrôles liés aux systèmes d'information pour
déceler les écarts par rapport à des critères, définis selon le type de
mission (audit financier, de conformité ou de performance).

\medskip
\termecle{Contrôles généraux}\index{contrôles généraux}~: procédures
manuelles ou automatisées qui visent à assurer la confidentialité,
l'intégrité et la disponibilité de l'information dans l'environnement où
les systèmes sont développés, gérés et exploités.

\medskip
\termecle{Contrôles d'application}\index{contrôles d'application}~:
procédures, dépendantes des TI, intégrées à un système et qui agissent sur
le traitement des transactions~: validation des données d'entrée,
exactitude du traitement, contrôle des sorties, intégrité des données de
référence.
\end{boxdef}

\subsection{La démarche du GUID 5100}

Le GUID~5100 suit les quatre étapes de tout audit public~: planification,
réalisation, rapport et suivi. Trois de ses apports sont particulièrement
utiles.

\paragraph{Les trois risques d'audit.} L'auditeur apprécie le
\emph{risque inhérent} (la probabilité que les caractéristiques mêmes d'un
système nuisent à la fonction qu'il remplit), le \emph{risque de contrôle}
(la probabilité que les contrôles mis en place ne suffisent pas) et le
\emph{risque de détection} (la probabilité que l'auditeur lui-même ne
décèle pas une défaillance). Le portail de paiement de l'ANPP, ouvert à
tous les fournisseurs de l'État, porte un risque inhérent élevé de
saturation.

\paragraph{Treize domaines.} Le GUID~5100 propose de délimiter la
vérification parmi treize domaines, dont la plupart sont adaptés de
l'ISO/IEC~27001~: politique de l'organisation en matière de TI,
structures de gouvernance, contrôles généraux, gestion des actifs,
développement et maintenance, exploitation, environnement physique,
ressources humaines, communications, sécurité de l'information,
conformité réglementaire, continuité et reprise, contrôles
d'application.

\paragraph{Conception, mise en œuvre, efficacité.} Le GUID~5100 distingue
trois niveaux d'examen d'un contrôle. Pour la \emph{conception}, un
entretien ou l'examen des règles documentées peut suffire. Pour la
\emph{mise en œuvre}, l'entretien ne suffit plus~: il faut une revue
structurée ou une analyse de données. Pour l'\emph{efficacité
opérationnelle}, il faut tester un échantillon d'opérations sur toute la
période.

\begin{boxprocedure}
\textbf{Recevoir une extraction de données.} Le GUID~5100 recommande que
toute extraction de données remise par l'entité auditée soit accompagnée
d'une lettre qui précise~:
\begin{itemize}
    \item la source des données et leur horodatage, avec l'empreinte
          numérique (valeur de hachage) du fichier, pour en garantir
          l'intégrité~;
    \item les paramètres d'extraction (requêtes utilisées, états
          exécutés)~;
    \item l'environnement d'origine (production ou autre).
\end{itemize}
À défaut de lettre, l'auditeur documente lui-même ces informations.
\end{boxprocedure}

\begin{boxscenario}{Opération MEFOU. La lettre de mission de la Chambre des comptes}
\textbf{Mandat.} Audit de conformité et de performance des systèmes
d'information de l'ANPP, période 2023--2026.

\medskip
\textbf{Délimitation retenue} (parmi les treize domaines)~: structures de
gouvernance des TI, développement et acquisition (le schéma directeur de
14~milliards FCFA), contrôles généraux (accès et changements) et
conformité réglementaire.

\medskip
\textbf{Critères.} Le régime financier de l'État (loi n°~2018/012), la loi
n°~2010/012, les procédures internes de l'agence et, comme référentiel de
bonnes pratiques, COBIT~2019.

\medskip
\textbf{Équipe.} Le magistrat vérificateur, M.~NJOYA, est assisté d'un
auditeur des SI certifié, recruté sur contrat, conformément à la
possibilité qu'ouvre le GUID~5100 de recourir à des ressources externes
supervisées.
\end{boxscenario}

% ============================================================================
\section{Tester les Contrôles Généraux Informatiques}
\label{sec:c11-itgc}
% ============================================================================

\subsection{Trois familles de contrôles}

Les contrôles généraux informatiques (Chapitre~\ref{chap:gouv-entreprise})
conditionnent la fiabilité de toutes les applications. L'auditeur les
teste en priorité, parce que leur défaillance invalide la confiance que
l'on peut accorder aux contrôles automatisés des applications~: si
n'importe qui peut modifier un programme, le contrôle que ce programme
réalise ne vaut plus rien.

\begin{tableaucours}{L{2.6cm}L{4.6cm}L{5.0cm}}{%
    \tableauHeader{Famille} &
    \tableauHeader{Contrôles typiques} &
    \tableauHeader{Test typique}
}
Accès & Création et suppression des comptes, revue périodique des droits,
comptes à privilèges & Rapprocher la liste des comptes actifs de la liste
des agents présents \\
Changements & Autorisation, test, séparation entre développement et
production, mise en production tracée & Vérifier, sur un échantillon de
changements, l'autorisation préalable et le test \\
Exploitation & Sauvegardes et tests de restauration, ordonnancement des
traitements, gestion des incidents & Vérifier les journaux de sauvegarde
et la preuve d'une restauration réussie \\
\end{tableaucours}
\vspace{-0.8em}
\begin{center}{\small\itshape Tableau~11.3. Les trois familles de contrôles
généraux informatiques}\end{center}

\subsection{L'échantillonnage}

On ne teste pas tous les comptes ni tous les changements~: on
échantillonne. La taille de l'échantillon dépend de la fréquence du
contrôle et du risque. Les cabinets d'audit utilisent couramment des
ordres de grandeur comme ceux du tableau suivant~; ce ne sont pas des
exigences normatives, et chaque organisation d'audit fixe sa propre
méthode.

\begin{tableaucours}{L{4.6cm}L{4.6cm}}{%
    \tableauHeader{Fréquence du contrôle} &
    \tableauHeader{Échantillon indicatif}
}
Annuelle & 1 \\
Trimestrielle & 2 \\
Mensuelle & 2 à 5 \\
Hebdomadaire & 5 à 15 \\
Quotidienne & 20 à 40 \\
Plusieurs fois par jour & 25 à 60 \\
\end{tableaucours}
\vspace{-0.8em}
\begin{center}{\small\itshape Tableau~11.4. Ordres de grandeur
d'échantillonnage pratiqués en audit (indicatifs)}\end{center}

\begin{boxscenario}{Opération MEFOU. Le test des accès}
\textbf{Test.} L'auditeur rapproche la liste des 312~comptes actifs du
système de paiement de l'ANPP et la liste du personnel au 30~juin~2026
fournie par les ressources humaines. Il examine en outre un échantillon de
25~comptes pour vérifier que les droits correspondent aux fonctions.

\medskip
\textbf{Résultats.} 14~comptes appartiennent à des agents partis de
l'agence, dont 3 depuis plus d'un an~; l'un de ces comptes a été utilisé
en mai~2026. Sur les 25~comptes examinés, 6 cumulent les droits de saisie
et de validation d'un paiement. Aucune revue des droits n'a été
documentée depuis 2023.

\medskip
\textbf{Qualification.} Le contrôle « revue périodique des accès » n'est
pas mis en œuvre~; le contrôle « séparation des tâches » est défaillant.
Ces écarts relèvent des objectifs DSS05 et DSS06 de COBIT et de
l'obligation de dispositifs de sécurité de l'article~26 de la loi
n°~2010/012 (Chapitre~\ref{chap:droit}).
\end{boxscenario}

\begin{boxmetier}{Auditeur SI}
\textbf{Tester un contrôle en cinq temps}
\begin{enumerate}
    \item Décrire le contrôle tel qu'il est censé fonctionner (qui, quoi,
          quand, avec quelle preuve).
    \item Vérifier sa conception~: le contrôle, s'il fonctionne, traite-t-il
          le risque~?
    \item Vérifier sa mise en œuvre sur un cas (test de cheminement).
    \item Tester son efficacité sur un échantillon couvrant la période.
    \item Documenter chaque pièce examinée, et conclure~: efficace,
          partiellement efficace, inefficace.
\end{enumerate}
\end{boxmetier}

% ============================================================================
\section{Rédiger un Constat d'Audit}
\label{sec:c11-constat}
% ============================================================================

\subsection{Les cinq éléments}

Un constat d'audit n'est pas une opinion~: c'est un raisonnement en cinq
éléments, que l'on retrouve dans les normes de l'IIA comme dans celles de
l'INTOSAI.

\begin{tableaucours}{L{2.8cm}L{4.6cm}L{4.8cm}}{%
    \tableauHeader{Élément} &
    \tableauHeader{Question} &
    \tableauHeader{Exemple (ANPP)}
}
Constatation & Que se passe-t-il réellement~? & 14 comptes d'agents partis
restent actifs~; l'un a été utilisé en mai~2026 \\
Critère & Que devrait-il se passer~? & Procédure interne~: désactivation
sous 48~heures~; loi n°~2010/012, art.~26 \\
Cause & Pourquoi l'écart existe-t-il~? & Les départs ne sont pas notifiés
par les RH à la division informatique \\
Conséquence & Quel risque ou quel effet~? & Paiement frauduleux possible
au nom d'un agent parti~; exposition juridique \\
Recommandation & Que faut-il faire~? & Notification automatique des
départs, revue trimestrielle des accès \\
\end{tableaucours}
\vspace{-0.8em}
\begin{center}{\small\itshape Tableau~11.5. Les cinq éléments d'un constat
d'audit}\end{center}

\begin{boxalert}
\textbf{Violation typique.} « La sécurité des accès est insuffisante et
doit être renforcée. » Cette phrase ne contient ni constatation chiffrée,
ni critère, ni cause~: elle ne peut être ni vérifiée ni mise en œuvre, et
l'audité peut la contester sans peine.

\smallskip
\textbf{Règle~:} chaque constat chiffre la constatation, cite son critère,
identifie la cause, mesure la conséquence, et formule une recommandation
qui traite la \emph{cause}, et non le symptôme.
\end{boxalert}

\subsection{Le suivi des recommandations}

Une recommandation non suivie est une assurance perdue. Le suivi relève de
la première ligne (qui met en œuvre), de l'audit interne (qui vérifie) et
de l'organe de gouvernance (qui exige). C'est l'objectif MEA02 de COBIT.

\begin{boxprocedure}
\textbf{Le tableau de suivi des recommandations.} Pour chaque
recommandation, le tableau présenté au comité d'audit indique~:
\begin{itemize}
    \item la date du rapport et la priorité de la recommandation~;
    \item le responsable de la mise en œuvre et l'échéance convenue~;
    \item l'état (non commencée, en cours, mise en œuvre, vérifiée)~;
    \item pour l'état « mise en œuvre », la pièce qui le prouve~;
    \item les reports d'échéance, avec leur motif et leur approbation.
\end{itemize}
La mention « en cours » sans échéance ni pièce, qui figurait dans le
rapport de la BANQUE SANAGA d'octobre~2025, est précisément ce qu'il faut
proscrire.
\end{boxprocedure}

% ============================================================================
\section{La Déontologie}
\label{sec:c11-deontologie}
% ============================================================================

\subsection{L'auditeur}

Le Code de déontologie de l'IIA (édition 2017) pose quatre principes, que
les Normes mondiales de 2024 reprennent dans leur domaine~II en y ajoutant
la conscience professionnelle~\cite{iia2017cripp}.

\begin{tableaucours}{L{2.8cm}L{9.4cm}}{%
    \tableauHeader{Principe} &
    \tableauHeader{Exigence (résumé)}
}
Intégrité & Elle fonde la confiance accordée au jugement de l'auditeur~:
honnêteté, diligence, responsabilité. \\
Objectivité & L'auditeur évalue équitablement tous les éléments
pertinents, sans se laisser influencer par ses intérêts ou par autrui. \\
Confidentialité & Il respecte la valeur et la propriété des informations
reçues, et ne les divulgue qu'avec autorisation, sauf obligation légale
ou professionnelle. \\
Compétence & Il met en œuvre les connaissances, les savoir-faire et
l'expérience que requiert sa mission. \\
\end{tableaucours}
\vspace{-0.8em}
\begin{center}{\small\itshape Tableau~11.6. Les principes du Code de
déontologie de l'IIA (résumé d'après l'édition française de 2017)}\end{center}

\subsection{Le DSI}

Le DSI n'est pas soumis à un code de déontologie professionnel
obligatoire, mais sa position l'expose à des situations délicates~: il
choisit des fournisseurs et dispose d'accès étendus aux données. Trois
règles de conduite en découlent.
\begin{itemize}
    \item \textbf{Déclarer ses conflits d'intérêts}, notamment dans le choix
          des prestataires, et s'abstenir dans les décisions concernées.
    \item \textbf{Rendre compte loyalement}, y compris des mauvaises
          nouvelles~: la note de la DSI de la BANQUE SANAGA du
          20~février~2026 en est l'exemple.
    \item \textbf{Ne pas faire obstacle à l'audit}~: donner accès aux pièces
          et aux systèmes, et ne pas « préparer » les données remises.
\end{itemize}

\begin{boiteRemarque}{L'auditeur peut-il conseiller~?}
Oui, et les normes de l'IIA le prévoient expressément~: l'audit interne
donne une assurance \emph{et} des conseils. Mais il y a une limite. Un
auditeur qui conçoit un dispositif ne peut pas, ensuite, en donner une
assurance objective. Dans une petite organisation, où l'audit interne
compte une ou deux personnes, cette tension est permanente. La solution
consiste à déclarer la situation au comité d'audit, et à faire réaliser
l'évaluation du dispositif par un tiers.
\end{boiteRemarque}

% ============================================================================
\separateur
% ============================================================================

\begin{boxresume}
\textbf{Ce qu'il faut retenir de ce chapitre~:}
\begin{itemize}
    \item \textbf{Trois lignes}~: le management porte le risque, les
          fonctions de deuxième ligne le surveillent, l'audit interne donne
          une assurance indépendante, sous la supervision de l'organe de
          gouvernance.
    \item \textbf{Les Normes mondiales de l'IIA (2024)} remplacent le cadre
          de 2017~: cinq domaines, dont la gouvernance de la fonction d'audit
          par le conseil.
    \item \textbf{Le GUID~5100 de l'INTOSAI} guide l'audit public des SI~:
          trois risques d'audit, treize domaines, et la distinction entre
          conception, mise en œuvre et efficacité des contrôles.
    \item \textbf{Les contrôles généraux informatiques} (accès,
          changements, exploitation) se testent en priorité, par
          échantillonnage.
    \item \textbf{Un constat a cinq éléments}~: constatation, critère,
          cause, conséquence, recommandation. La recommandation traite la
          cause.
    \item \textbf{Le suivi est une responsabilité de gouvernance}~: une
          recommandation « en cours » sans échéance ni pièce est une
          assurance perdue.
    \item \textbf{La déontologie} repose sur l'intégrité, l'objectivité, la
          confidentialité et la compétence~; elle concerne aussi le DSI.
\end{itemize}
\end{boxresume}

\bigskip

\begin{boxexo}[title={\ding{45}\quad Exercice~11.1. Quelle ligne~? \niveaubase}]
Rattachez chacun des acteurs suivants à la ligne qui convient (première,
deuxième, troisième, ou prestataire d'assurance externe), et justifiez.
\begin{enumerate}[(a)]
    \item Le responsable de l'exploitation de la BANQUE SANAGA.
    \item Le RSSI, rattaché au directeur des risques.
    \item Le commissaire aux comptes.
    \item La directrice de l'audit interne.
    \item Le cabinet agréé qui réalise l'audit de sécurité prévu par la loi
          n°~2010/012.
\end{enumerate}
\end{boxexo}

\begin{boxexo}[title={\ding{45}\quad Exercice~11.2. Rédiger un constat \niveaumoyen}]
À partir des résultats du test des accès de l'ANPP (section~\ref{sec:c11-itgc}),
rédigez le constat relatif à la séparation des tâches (6~comptes sur
25 cumulant saisie et validation d'un paiement).
\begin{enumerate}[(a)]
    \item Rédigez les cinq éléments du constat.
    \item Proposez une priorité (élevée, moyenne, faible) et justifiez.
    \item Proposez l'indicateur qui permettra au comité d'audit de suivre la
          mise en œuvre.
\end{enumerate}
\end{boxexo}

\begin{boxexo}[title={\ding{45}\quad Exercice~11.3. Planifier l'audit des changements de la BANQUE SANAGA \niveauexpert}]
Le comité d'audit demande à Mme~EKAMBI un audit de la gestion des
changements après la panne.
\begin{enumerate}[(a)]
    \item Rédigez l'objectif, le périmètre et les critères de la mission.
    \item Listez les contrôles à tester, en distinguant conception, mise en
          œuvre et efficacité, et indiquez pour chacun la pièce attendue.
    \item Déterminez la taille de l'échantillon de changements à examiner,
          en justifiant au regard de la fréquence des changements.
    \item Mme~EKAMBI avait elle-même proposé, en 2025, la procédure de
          gestion des changements. Quelle difficulté déontologique cela
          pose-t-il, et comment la traiter~?
\end{enumerate}
\end{boxexo}

\bigskip

\begin{boxinfo}
\textbf{Transition.} La Partie~\ref{part:normatif} a placé la gouvernance
des SI sous le regard des normes, de la loi et de l'audit. La
Partie~\ref{part:pratique} descend au niveau de la pratique quotidienne.
Le Chapitre~\ref{chap:itil-svs} ouvre la gestion des services avec ITIL~4
et son système de valeur des services.
\end{boxinfo}
